\chapter{Exercises with Solutions}
\label{app:exercises}

Each exercise certifies one specific step rather than a topic. Solutions are complete enough to be
checked line by line.

\begin{exercise}\label{ex1}
Fix $\omega\neq0$ and show that $I_{T}(\omega)=2\sin(\omega T)/\omega$ has no limit as
$T\to\infty$.
\end{exercise}

\begin{proof}[Solution]
Let $a=|\omega|$ and set $T_{n}=(\pi/2+2\pi n)/a$, $S_{n}=(3\pi/2+2\pi n)/a$. Then
$\sin(aT_{n})=1$ and $\sin(aS_{n})=-1$. Since $\sin(\omega T)=\operatorname{sgn}(\omega)\sin(aT)$,
\begin{equation}
I_{T_{n}}(\omega)=\operatorname{sgn}(\omega)\frac{2}{a},
\qquad
I_{S_{n}}(\omega)=-\operatorname{sgn}(\omega)\frac{2}{a},
\end{equation}
and both $T_{n},S_{n}\to\infty$. Two subsequences with distinct limits preclude convergence.
\end{proof}

\begin{exercise}\label{ex2}
Compute $\int_{-a}^{a}G_{\ep}(\omega)\dd\omega$ and explain why a fixed interval eventually contains
almost all the mass.
\end{exercise}

\begin{proof}[Solution]
\begin{equation}
\int_{-a}^{a}\frac{2\ep}{\ep^{2}+\omega^{2}}\dd\omega
=2\left[\arctan\frac{\omega}{\ep}\right]_{-a}^{a}
=4\arctan\frac{a}{\ep}\longrightarrow4\cdot\frac{\pi}{2}=2\pi .
\end{equation}
The order of quantifiers is the point: $a$ is fixed and $\ep\downarrow0$, so $\arctan(a/\ep)\to\pi/2$.
Sending $a\to\infty$ first would instead compute the total mass of each fixed member, which is a
different statement.
\end{proof}

\begin{exercise}\label{ex3}
For $\varphi(x)=(1+x)e^{-x^{2}}$ compute $\ip{\delta}{\varphi}$, $\ip{\delta'}{\varphi}$,
$\ip{\delta''}{\varphi}$.
\end{exercise}

\begin{proof}[Solution]
$\varphi(0)=1$. Differentiating,
$\varphi'(x)=e^{-x^{2}}\big(1-2x-2x^{2}\big)$, so $\varphi'(0)=1$. A second differentiation at $0$
gives $\varphi''(0)=-2$. Hence, by \eqref{eq:deltan},
\begin{equation}
\ip{\delta}{\varphi}=1,\qquad
\ip{\delta'}{\varphi}=-1,\qquad
\ip{\delta''}{\varphi}=-2 .
\end{equation}
\end{proof}

\begin{exercise}\label{ex4}
Compute $\F[3-2t+t^{2}]$ and verify the symmetry statement of
Proposition~\ref{prop:symmetries} in distributional form.
\end{exercise}

\begin{proof}[Solution]
By Theorem~\ref{thm:tn} and linearity,
\begin{equation}
\F[3-2t+t^{2}]=6\pi\delta-4\pi i\,\delta'-2\pi\,\delta'' .
\end{equation}
The function $3+t^{2}$ is real and even; it contributes $6\pi\delta-2\pi\delta''$, which is real and
even as a distribution. The term $-2t$ is real and odd and contributes $-4\pi i\delta'$, which is
purely imaginary and odd. Since $\delta^{(n)}$ is even for $n$ even and odd for $n$ odd
(it reads off $\varphi^{(n)}(0)$, and reflection changes $n$-th derivatives by $(-1)^{n}$), the
elementary symmetries are reproduced exactly.
\end{proof}

\begin{exercise}\label{ex5}
Prove $x\delta'=-\delta$ and $x\,\PV\frac{1}{x}=1$, and deduce $i\omega\hf=1$ for
$\hf=\pi\delta-i\PV\frac{1}{\omega}$.
\end{exercise}

\begin{proof}[Solution]
For the first, $\ip{x\delta'}{\varphi}=\ip{\delta'}{x\varphi}=-(x\varphi)'(0)=-\varphi(0)$.
For the second, by \eqref{eq:PVint} with $x\varphi$ in place of $\varphi$,
\begin{equation}
\left\langle x\PV\tfrac{1}{x},\varphi\right\rangle
=\int_{0}^{\infty}\frac{x\varphi(x)+x\varphi(-x)}{x}\dd x
=\int_{0}^{\infty}\big(\varphi(x)+\varphi(-x)\big)\dd x=\int_{\R}\varphi .
\end{equation}
Applying multiplication by $i\omega$ to the candidate spectrum,
\begin{equation}
i\omega\Big[\pi\delta-i\PV\tfrac{1}{\omega}\Big]
=i\pi\underbrace{\omega\delta}_{=0}+\underbrace{\omega\PV\tfrac{1}{\omega}}_{=1}=1 ,
\end{equation}
which is \eqref{eq:iomegaH}. This also displays the ambiguity: adding $c\delta$ changes nothing.
\end{proof}

\begin{exercise}\label{ex6}
Show that the value assigned to $H(0)$ does not affect the distribution it defines.
\end{exercise}

\begin{proof}[Solution]
Let $H_{1},H_{2}$ differ only at $t=0$, with $|H_{1}(0)-H_{2}(0)|=c<\infty$. For
$\varphi\in C_{c}^{\infty}$,
\begin{equation}
\ip{H_{1}}{\varphi}-\ip{H_{2}}{\varphi}
=\int_{\R}(H_{1}-H_{2})\varphi\dd t=0 ,
\end{equation}
because the integrand vanishes except at a single point and the Lebesgue integral ignores null sets.
\end{proof}

\begin{exercise}\label{ex7}
Explain why no single $\Lone$ function dominates the family $K_{\ep}$, and why the substitution
$x=\ep u$ repairs the situation.
\end{exercise}

\begin{proof}[Solution]
If $K_{\ep}\le g\in\Lone$ for all $\ep$, then $K_{\ep}\to0$ pointwise off the origin and dominated
convergence gives $\int K_{\ep}\to0$, contradicting $\int K_{\ep}=1$.

The substitution does not produce a dominator for $K_{\ep}$; it changes the object being integrated:
\begin{equation}
\int_{\R}K_{\ep}(x)\varphi(x)\dd x=\frac{1}{\pi}\int_{\R}\frac{\varphi(\ep u)}{1+u^{2}}\dd u ,
\end{equation}
and now the integrand is dominated by $\|\varphi\|_{\infty}/(\pi(1+u^{2}))$, independent of $\ep$ and
integrable. The varying kernel has been replaced by a fixed weight, at the cost of letting the test
function depend on $\ep$.
\end{proof}

\begin{exercise}\label{ex8}
Exhibit a family with unit mass whose distributional limit is $0$, and reconcile this with mass
conservation.
\end{exercise}

\begin{proof}[Solution]
Take $k_{n}=\tfrac12\mathbf{1}_{[n-1,n+1]}$. For $\varphi\in\Sw$,
$|\ip{k_{n}}{\varphi}|\le\tfrac12\int_{n-1}^{n+1}|\varphi|=O(n^{-2})\to0$. There is no contradiction:
$\int k_{n}=1$ is a statement about the functional $\varphi\equiv1$, and $1\notin\Sw$. Total mass is
therefore not a continuous functional on $\Sw'$; the mass is conserved but escapes to infinity, where
no rapidly decaying probe can detect it.
\end{proof}

\begin{exercise}\label{ex9}
Using only the dual definition, prove $\F[\delta]=1$ and $\F[1]=2\pi\delta$, and identify where
inversion enters.
\end{exercise}

\begin{proof}[Solution]
\begin{equation}
\ip{\F\delta}{\varphi}=\ip{\delta}{\ph}=\ph(0)=\int_{\R}\varphi=\ip{1}{\varphi},
\end{equation}
because $\ph(0)=\int\varphi(t)\dd t$. For the second,
\begin{equation}
\ip{\F1}{\varphi}=\ip{1}{\ph}=\int_{\R}\ph(t)\dd t=2\pi\varphi(0)=\ip{2\pi\delta}{\varphi},
\end{equation}
and the middle equality is exactly \eqref{eq:inversion} evaluated at $t=0$.
\end{proof}

\begin{exercise}\label{ex10}
Compute $\F[H(\cdot-a)]$ and explain why the $\delta$ term is unaffected by the shift.
\end{exercise}

\begin{proof}[Solution]
By \eqref{eq:timeshift},
\begin{equation}
\F[H(\cdot-a)](\omega)=e^{-ia\omega}\Big[\pi\delta(\omega)-i\PV\frac{1}{\omega}\Big]
=\pi\delta(\omega)-i\,e^{-ia\omega}\PV\frac{1}{\omega}.
\end{equation}
The factor $e^{-ia\omega}$ is a smooth multiplier and $e^{-ia\cdot0}=1$, so by \eqref{eq:adelta}
\begin{equation}
e^{-ia\omega}\,\delta(\omega)=\delta(\omega).
\end{equation}
The principal value term keeps the factor, because $\PV(1/\omega)$ is not concentrated at a point.
\end{proof}

\begin{exercise}\label{ex11}
Locate the two sign changes that produce the $+i$ in $\F(tT)=i\,\partial_{\omega}\widehat T$.
\end{exercise}

\begin{proof}[Solution]
First, from $\F[\varphi']=i\omega\ph(\omega)$ with the variable renamed,
$t\ph(t)=-i(\varphi')^{\wedge}(t)$; the factor $1/i=-i$ is the first sign. Second, the definition of
the distributional derivative gives
$\ip{\widehat T'}{\varphi}=-\ip{\widehat T}{\varphi'}$, a second minus. Then
\begin{equation}
\ip{\F(tT)}{\varphi}=\ip{T}{t\ph}
=-i\ip{T}{(\varphi')^{\wedge}}
=-i\ip{\widehat T}{\varphi'}
=i\ip{\widehat T'}{\varphi}.
\end{equation}
The two negatives cancel, leaving $+i$.
\end{proof}

\begin{exercise}\label{ex12}
For $K_{\ep}'(x)=-2\ep x/(\pi(x^{2}+\ep^{2})^{2})$ determine the location and size of the extrema,
and compute $\int|K_{\ep}'|$.
\end{exercise}

\begin{proof}[Solution]
Differentiating $\log|K_{\ep}'|$ gives the condition
$\frac{\dd}{\dd x}\big(\log x-2\log(x^{2}+\ep^{2})\big)=1/x-4x/(x^{2}+\ep^{2})=0$, that is
$\ep^{2}-3x^{2}=0$, so $x=\pm\ep/\sqrt3$. Substituting,
\begin{equation}
|K_{\ep}'(\ep/\sqrt3)|
=\frac{2\ep\cdot\ep/\sqrt3}{\pi\big(\ep^{2}+\ep^{2}/3\big)^{2}}
=\frac{2\ep^{2}/\sqrt3}{\pi\cdot 16\ep^{4}/9}
=\frac{9}{8\sqrt3\,\pi\ep^{2}} .
\end{equation}
For the total variation, use symmetry and the substitution $u=x^{2}+\ep^{2}$:
\begin{equation}
\int_{\R}|K_{\ep}'|\dd x
=2\int_{0}^{\infty}\frac{2\ep x}{\pi(x^{2}+\ep^{2})^{2}}\dd x
=\frac{4\ep}{\pi}\cdot\frac{1}{2\ep^{2}}
=\frac{2}{\pi\ep}.
\end{equation}
Since $\int K_{\ep}'=0$, the positive and negative masses are each $1/(\pi\ep)$, both diverging.
\end{proof}

\begin{exercise}\label{ex13}
Show that $e^{x^{2}}$ defines a distribution but not a tempered distribution.
\end{exercise}

\begin{proof}[Solution]
Local integrability is clear, so Theorem~\ref{thm:embedding} applies. Suppose it were tempered, so
that \eqref{eq:Sprime-est} held. Fix $0\not\equiv\psi\in C_{c}^{\infty}$ nonnegative and set
$\psi_{R}(x)=\psi(x-R)$. Then $p_{m,n}(\psi_{R})=O(R^{m})$ for each fixed $m,n$, so the right side of
\eqref{eq:Sprime-est} is $O(R^{M})$, whereas
\begin{equation}
\ip{e^{x^{2}}}{\psi_{R}}\ge e^{(R-r)^{2}}\int\psi
\end{equation}
grows faster than any polynomial. Contradiction.
\end{proof}

\begin{exercise}\label{ex14}
Prove the second identity of Lemma~\ref{lem:two-identities}, namely
$\omega^{m}\ph(\omega)=i^{-m}\F[\varphi^{(m)}](\omega)$.
\end{exercise}

\begin{proof}[Solution]
By Theorem~\ref{thm:derivt} applied to $\varphi^{(m)}$,
$\F[\varphi^{(m)}]=(i\omega)^{m}\ph(\omega)$. Multiplying both sides by $i^{-m}$ and using
$i^{-m}i^{m}=1$ gives the claim.
\end{proof}

\begin{exercise}\label{ex15}
Show that if $f$ is real and even then $\fh$ is real and even, and verify this for
$G_{\ep}=\F[e^{-\ep|\cdot|}]$.
\end{exercise}

\begin{proof}[Solution]
Evenness of $f$ gives, after $t\mapsto-t$,
\begin{equation}
\fh(-\omega)=\int f(t)e^{i\omega t}\dd t=\int f(-s)e^{-i\omega s}\dd s=\int f(s)e^{-i\omega s}\dd s=\fh(\omega),
\end{equation}
and $f$ real gives $\overline{\fh(\omega)}=\int f(t)e^{i\omega t}\dd t=\fh(-\omega)$, hence
$\fh=\overline{\fh}$. For $f_{\ep}=e^{-\ep|t|}$, which is real and even, one gets
$G_{\ep}=2\ep/(\ep^{2}+\omega^{2})$, visibly real and even.
\end{proof}

\begin{exercise}\label{ex16}
Where exactly does the proof of Theorem~\ref{thm:inversion} use Fubini, and what makes the double
integral absolutely convergent?
\end{exercise}

\begin{proof}[Solution]
Fubini is used to pass from
$\frac{1}{2\pi}\iint\varphi(s)e^{-i\omega s}e^{i\omega t}e^{-\eta\omega^{2}/2}\dd s\dd\omega$ to the
iterated integral with the inner integral over $\omega$. Absolute convergence holds because the modulus
of the integrand factors as $|\varphi(s)|\,e^{-\eta\omega^{2}/2}$, and the product of two integrable
functions is integrable on the product space: $\varphi\in\Lone$ since $\varphi\in\Sw$, and
$e^{-\eta\omega^{2}/2}$ is integrable in $\omega$. The Gaussian factor is indispensable; without it the
double integral is not absolutely convergent and the interchange is unjustified.
\end{proof}
